% The_Representational_Regress.tex — Full paper
% Based on the proved Lean artifact (0 sorry, 0 custom axioms)
% Companion to representational-regress-lean

\documentclass[11pt]{article}
\usepackage[margin=1in]{geometry}
\usepackage{amsmath,amssymb,amsthm}
\usepackage{hyperref}
\usepackage{enumitem}
\usepackage{booktabs}
\usepackage{xcolor}
\usepackage{stmaryrd}

\newtheorem{theorem}{Theorem}[section]
\newtheorem{lemma}[theorem]{Lemma}
\newtheorem{proposition}[theorem]{Proposition}
\newtheorem{corollary}[theorem]{Corollary}
\newtheorem{definition}[theorem]{Definition}
\newtheorem{remark}[theorem]{Remark}

\newcommand{\leanref}[1]{\texttt{\small #1}}
\newcommand{\novel}{\textcolor{blue!70!black}{[Novel]}}
\newcommand{\known}{\textcolor{gray}{[Known math]}}

% Title: mathematical frame first; consciousness and other readings are instantiations (§7).
\title{The Representational Regress:\\
  Machine-Checked Topology, Diagonal Arguments,\\
  and Boundary Obstructions for Representational Systems}
\author{[Authors]}
\date{2026}

\begin{document}
\maketitle

% ============================================================================
% ABSTRACT
% ============================================================================
\begin{abstract}
We present a Lean~4 formalization, verified against Mathlib with zero
\texttt{sorry} and zero custom axioms, for a family of claims about
\emph{representational systems}---categories with a non-trivial self-representation
endomorphism---and about \emph{intrinsic} planar-chartability as a topological
sensor on compact models.

\begin{enumerate}[label=(\roman*)]
\item \textbf{Categorical strand.} A representational system generates infinitely
  many distinct meta-levels; Lawvere's fixed-point theorem, packaged in
  \texttt{Type} and \texttt{MonoidalClosed(Type)}, does not dissolve an
  object/morphism distinction formalized by the inductive type
  \texttt{OntologicalSlot} (the ``wall'' remains).
\item \textbf{Topological strand.} The closed cylinder $S^1 \times [0,1]$ and the
  M\"obius strip (quotient of $[0,1]^2$ by $(0,t)\sim(1,1-t)$) are \emph{not}
  homeomorphic. The engine is \texttt{ChartableR2} (open neighborhood
  $\cong\mathbb{R}^2$): boundary points are exactly non-\texttt{ChartableR2} on
  both models, and the boundary \emph{subspaces} are incompatible (connected vs.\
  disconnected). A \textbf{generic} lemma
  (\texttt{not\_homeomorphic\_of\_chartableR2\_boundary\_contrast}) abstracts this
  pattern for any future pair of spaces with the same sensor hypotheses.
\item \textbf{Paper-facing instantiations (mostly re-exports).} The library adds
  named packaging for \emph{symbol grounding} (\texttt{SymbolSystem}),
  Lawvere-style \emph{no universal self-interpreter} lemmas
  (\texttt{LawvereRegressBridge}), bundled \texttt{ChartableR2BoundaryModel}s for
  the cylinder and M\"obius strip, and a \textbf{hypothesis-only} schematic
  \texttt{QuantumObserverChainHypothesis} recording what a future formal
  observer/measurement-chain model would need to reuse the same obstruction---%
  \emph{no} physics axioms and no quantum formalism in the artifact.
\end{enumerate}

\medskip
\noindent\textbf{Consciousness as one instantiation.}
The mathematics does not mention consciousness. A \emph{philosophical} premise
P2 (that a form of self-presence is modeled by M\"obius-type rather than
cylinder-type topology) yields a \emph{conditional} gloss: no representational
system whose world bears the cylinder-side invariant can be globally
homeomorphic to that target. The distinction is invariant under scaling,
iteration, and continuous deformation of the representational layer.

\medskip
The artifact comprises ${\sim}25$ Lean modules; we summarize the
\texttt{OntologicalSlot} construction, the seam--equator chart technique, and
honest limits in the text and appendices.
\end{abstract}

% ============================================================================
% 1. INTRODUCTION
% ============================================================================
\section{Introduction}
\label{sec:intro}

% ~3-4 pages. Structure:

\subsection{The Problem}
Representational systems---symbolic AI, formal theories, any setting where a
structure represents itself by iterating a fixed construction---support arbitrarily
deep \emph{meta}-levels. Classical diagonal and fixed-point arguments constrain
how much ``self-capture'' such a layer can achieve. Independently, one may ask
whether intrinsic topology distinguishes different \emph{kinds} of ``inside vs.\
outside'' organization when a global homeomorphism is the notion of sameness.

This paper reports a single Lean artifact that makes a substantial fragment of
that story adversarially checkable. We state the mathematics without
phenomenological vocabulary; \S\ref{sec:instantiations} explains how
consciousness, grounding, and related readings attach as \emph{optional}
instantiations.

\subsection{Summary of Results}
The artifact proves, with zero \texttt{sorry} and zero custom axioms:
\begin{enumerate}[label=(\alph*)]
\item \textbf{Categorical regress and Lawvere wall} (\S\ref{sec:regress}):
  unbounded distinct iterates of \texttt{represent}; Lawvere in \texttt{Type} /
  CCC flavor; \texttt{OntologicalSlot} disjointness so the wall does not
  dissolve into a bare object witness.
\item \textbf{Generic boundary obstruction} (\S\ref{sec:generic},
  \S\ref{sec:mfinal}): if two compact models carry \texttt{ChartableR2}
  boundary sensors and their boundary subspaces are not homeomorphic, the
  ambient spaces are not homeomorphic. The cylinder vs.\ M\"obius strip is the
  flagship witness (\textbf{M-FINAL}).
\item \textbf{Extended master bundle} (\S\ref{sec:extended}):
  \texttt{representational\_regress\_master\_extended} adjoins the $1$D
  half-space obstruction and M-FINAL to the core master claim.
\item \textbf{Named instantiations in Lean} (\S\ref{sec:instantiations}):
  \texttt{SymbolSystem}, \texttt{LawvereRegressBridge},
  \texttt{ChartableR2BoundaryModel} instances, and a hypothesis-only
  \texttt{QuantumObserverChainHypothesis} schematic.
\end{enumerate}

\subsection{Contributions and Novelty Claims}
\begin{enumerate}[label=(\arabic*)]
\item \textbf{\texttt{OntologicalSlot}} --- machine-checked tagging: morphism
  slots vs.\ object slots; ``wall'' lemmas are definitional but mathematically
  load-bearing for the paper's reading of Lawvere-style tallies. [\novel]
\item \textbf{\texttt{ChartableR2} spine + seam--equator charts} --- intrinsic
  $\mathbb{R}^2$ chartability as boundary sensor; piecewise unfolding at the
  quotient equator. [\novel\ application / technique]
\item \textbf{Generic obstruction lemma}
  (\texttt{not\_homeomorphic\_of\_chartableR2\_boundary\_contrast}) ---
  reusable template beyond the cylinder/M\"obius pair. [\novel]
\item \textbf{Instantation-facing API} --- \texttt{SymbolSystem},
  \texttt{LawvereRegressBridge}, concrete boundary models, QM schematic;
  supports multidisciplinary citation without duplicating proofs. [\novel]
\item \textbf{Not novel:} Lawvere (1969), standard surface topology, Cantorian
  regress; our contribution is verification, packaging, and application shape.
\end{enumerate}

\subsection{Paper Organization}
Section~\ref{sec:regress}: representational systems, regress, Lawvere,
\texttt{OntologicalSlot}, core master theorem.
Section~\ref{sec:models}--\ref{sec:mfinal}: cylinder and M\"obius models,
\texttt{ChartableR2}, generic lemma, seam charts, M-FINAL.
Section~\ref{sec:extended}: extended master and instantiations (consciousness,
grounding, diagonal arguments, QM schematic).
Sections~\ref{sec:limits}--\ref{sec:conclusion}: limits, related work, conclusion.
Appendices: Lean module map, theorem index, reproduction.

% ============================================================================
% 2. CATEGORICAL STRAND: THE REPRESENTATIONAL REGRESS
% ============================================================================
\section{The Representational Regress}
\label{sec:regress}

% ~4 pages.

\subsection{The \texttt{RepresentationalSystem} Structure}
\label{sec:repsys}

\begin{definition}[\leanref{Basic.lean}: \leanref{RepresentationalSystem}]
A \emph{representational system} is a tuple $(C, A, \mathrm{represent}, \mathrm{iter\_injective})$ where:
\begin{itemize}
\item $C$ is a category,
\item $A$ is an object of $C$,
\item $\mathrm{represent} : A \to A$ is an endomorphism,
\item $\mathrm{iter\_injective}$: for all $n \neq m$, $\mathrm{represent}^n \neq \mathrm{represent}^m$ in $\mathrm{End}(A)$.
\end{itemize}
\end{definition}

% Discuss: why bundle iter_injective (captures "non-trivial self-reference");
% metaRegressArrow, metaOver, metaRepresent as the indexed tower;
% slice-object packaging via Over.mk.
%
% Lean refs: representIter, metaRegressArrow, metaOver, Over_mk_inj_parallel

\subsection{Non-Termination and Unboundedness}

\begin{theorem}[\known\ \leanref{Regress.lean}]
\label{thm:regress}
For any representational system $R$:
\begin{enumerate}[label=(\alph*)]
\item (\leanref{regress\_no\_termination}) Distinct indices $n \neq m$ yield distinct meta-level arrows.
\item (\leanref{regress\_iterates\_unbounded}) For every bound $B$, there exists a level $L > B$ with a distinct iterate.
\item (\leanref{regress\_is\_infinite}, \leanref{meta\_range\_infinite}) The set of meta-levels is infinite.
\end{enumerate}
\end{theorem}

% Proof sketch: direct from iter_injective + monoid exponentiation.
% Novelty: zero (standard); contribution is machine-checking.

\subsection{The Lawvere Wall}
\label{sec:lawvere}

\begin{theorem}[\known\ \leanref{Lawvere.lean}]
\label{thm:lawvere}
(\leanref{lawvere\_fixed\_point\_Type}) If $s : A \to (A \to B)$ is surjective, then every $f : B \to B$ has a fixed point.
\end{theorem}

\begin{theorem}[\known\ \leanref{LawvereCCCType.lean}]
\label{thm:lawvere-ccc}
(\leanref{lawvere\_fixed\_point\_MonoidalClosedType}) The same conclusion holds when the universal-enumerator hypothesis is recast as surjectivity of $\mathrm{curry}(\mathrm{lawvereBinary}\ s)$ in $\mathrm{MonoidalClosed}(\mathrm{Type}\ u)$.
\end{theorem}

% Discuss: Lawvere 1969, diagonal argument, why we formalize both flavors.
% Lean refs: lawvereBinary, curry_lawvereBinary, lawvere_universal_iff_surjective_curry.

\subsection{The \texttt{OntologicalSlot} Construction}
\label{sec:ontoslot}

\begin{definition}[\novel\ \leanref{Basic.lean}: \leanref{OntologicalSlot}]
\label{def:ontoslot}
For a category $C$ with terminal object and an object $A$:
\[
\mathrm{OntologicalSlot}(C, A) \;:=\;
  \mathrm{endo}(f : A \to A) \;\mid\; \mathrm{obj}(g : \top_C \to A)
\]
an inductive type with two constructors, disjoint by definition.
\end{definition}

\begin{theorem}[\novel\ \leanref{FixedPoints.lean}]
\label{thm:wall}
(\leanref{lawvere\_wall\_is\_not\_dissolution}) For any representational system $R$ and any global element $\mathrm{fp} : \top_C \to A$,
\[
  \mathrm{OntologicalSlot.endo}(R.\mathrm{represent}) \;\neq\; \mathrm{OntologicalSlot.obj}(\mathrm{fp}).
\]
\end{theorem}

% The proof is trivially constructor disjointness — and that is the point.
% Discuss: the philosophical content IS the triviality. Lawvere's fixed point
% is still an object, not a dissolution. Contrast with Hofstadter's "strange
% loops" (which remain within the orientable topology).
%
% Additional theorems:
%   OntologicalSlot.endo_injective
%   representational_tower_preserves_slot (distinct n => distinct endo slots)
%   no_slot_collapse (no endo slot equals an obj slot)

\subsection{The Core Master Theorem}

\begin{theorem}[\leanref{Main.lean}]
\label{thm:master}
(\leanref{representational\_regress\_master}) For any representational system $R$:
\[
  \mathrm{representational\_regress\_master\_claim}(R)
\]
i.e., the conjunction of unbounded regress, Lawvere wall non-dissolution, and $T_2$ invariance under homeomorphism.
\end{theorem}

% This is the SPEC_001 anchor; topology punchline is added in §7.

% ============================================================================
% 3. TOPOLOGICAL MODELS
% ============================================================================
\section{Topological Models: Cylinder and M\"obius Strip}
\label{sec:models}

% ~3 pages.

\subsection{The Closed Cylinder}
\label{sec:cylinder}

\begin{definition}[\leanref{CylinderBoundary.lean}]
$\mathrm{ClosedCylinder} := S^1 \times [0,1]$ with the product topology and manifold structure $\mathfrak{R}_1 \times \mathfrak{R}_1^\partial$.
\end{definition}

\begin{theorem}[\known\ \leanref{CylinderBoundary.lean}]
\label{thm:cyl-boundary}
The manifold boundary of \emph{ClosedCylinder} decomposes as:
\[
  \partial(\mathrm{ClosedCylinder}) = F_{\mathrm{bot}} \sqcup F_{\mathrm{top}}
\]
where $F_{\mathrm{bot}} = S^1 \times \{0\}$ and $F_{\mathrm{top}} = S^1 \times \{1\}$ are each connected, and their union is \emph{disconnected}.
\end{theorem}

% Lean: closedCylinder_boundary_eq_union, isConnected_closedCylinderBotFace,
%   closedCylinder_boundary_faces_disjoint, not_isConnected_closedCylinderBoundaryUnion.
% Via Mathlib boundary_product.

\subsection{The M\"obius Strip}
\label{sec:mobius}

\begin{definition}[\leanref{CylinderMobius.lean}]
\label{def:mobius-rel}
The \emph{M\"obius equivalence} $\sim$ on $[0,1]^2$ is the equivalence relation generated by
$(0, t) \sim (1, 1-t)$ for $t \in [0,1]$, extended by reflexivity, symmetry, and transitivity.
The \emph{M\"obius strip} is $\mathrm{MobiusStrip} := [0,1]^2 / {\sim}$.
\end{definition}

\begin{theorem}[\leanref{CylinderMobius.lean}]
\label{thm:mobius-props}
\leavevmode
\begin{enumerate}[label=(\alph*)]
\item (\leanref{mobiusRel0Graph\_isClosed}) The graph of $\sim$ is closed in $[0,1]^2 \times [0,1]^2$.
\item (\leanref{CompactSpace MobiusStrip}) $\mathrm{MobiusStrip}$ is compact (quotient of compact).
\item (\leanref{instT2SpaceMobiusStrip}) $\mathrm{MobiusStrip}$ is Hausdorff (closed graph + compact domain).
\item (\leanref{isConnected\_mobiusStripBoundary}) The boundary band $\partial_{\mathrm{model}} := \pi(\{t = 0\} \cup \{t = 1\})$ is connected (one component).
\end{enumerate}
\end{theorem}

% Discuss: mobiusRel₀ as a direct equivalence (no transitive closure needed
% beyond the explicit three-case definition); the closed-graph method for T2;
% boundary connectivity through corner identification.

\subsection{Boundary Subspace Contrast}

\begin{theorem}[\leanref{MobiusCylinderBoundaryContrast.lean}]
\label{thm:boundary-contrast}
(\leanref{mobiusStripBoundary\_not\_homeomorphic\_closedCylinderBoundaryUnion})
\[
  \mathrm{IsEmpty}\!\left(\mathrm{Subtype}(\partial_{\mathrm{Mob}}) \simeq_t \mathrm{Subtype}(\partial_{\mathrm{Cyl}})\right).
\]
The M\"obius boundary band (connected) and the cylinder boundary union (disconnected) are not homeomorphic as subspaces.
\end{theorem}

% ============================================================================
% 4. THE CHARTABLER2 FRAMEWORK
% ============================================================================
\section{The \texttt{ChartableR2} Framework}
\label{sec:chartable}

% ~3 pages. This is the key technical innovation for the non-homeomorphism.

\subsection{Definition and Basic Properties}

\begin{definition}[\leanref{ChartableR2Neighbor.lean}]
\label{def:chartable}
A point $x$ of a topological space $X$ is \emph{$\mathrm{ChartableR2}$} if there exists an open set $U \ni x$ such that $U \cong_t \mathbb{R}^2$.
\end{definition}

\begin{theorem}[\leanref{ChartableR2Neighbor.lean}]
\label{thm:chartable-props}
\leavevmode
\begin{enumerate}[label=(\alph*)]
\item (\leanref{isOpen\_setOf\_chartableR2}) The set of \texttt{ChartableR2} points is open.
\item (\leanref{isClosed\_setOf\_not\_chartableR2}) Its complement is closed.
\item (\leanref{not\_chartableR2\_of\_tendsto\_seq\_atTop}) Sequential stability: limits of non-chartable points are non-chartable.
\end{enumerate}
\end{theorem}

\subsection{Invariance Under Homeomorphism}

\begin{theorem}[\leanref{ChartableR2Bridge.lean}]
\label{thm:chartable-invar}
(\leanref{Homeomorph.chartableR2\_iff}) If $h : X \simeq_t Y$ is a homeomorphism, then $\mathrm{ChartableR2}(x) \iff \mathrm{ChartableR2}(h(x))$.
\end{theorem}

\subsection{The Generic Boundary-Contrast Obstruction}
\label{sec:generic}

The following theorem is the \textbf{universal engine} behind M-FINAL and any
future surface pair separated by $\mathrm{ChartableR2}$-detected boundary topology.

\begin{theorem}[\novel\ \leanref{ChartableR2Bridge.lean}]
\label{thm:generic}
(\leanref{not\_homeomorphic\_of\_chartableR2\_boundary\_contrast})
Let $X$ and $Y$ be topological spaces with predicates $b_X \subseteq X$ and $b_Y \subseteq Y$ such that:
\begin{enumerate}[label=(\roman*)]
\item $\forall x,\; x \in b_X \iff \neg\,\mathrm{ChartableR2}(x)$,
\item $\forall y,\; y \in b_Y \iff \neg\,\mathrm{ChartableR2}(y)$,
\item $\mathrm{IsEmpty}(\mathrm{Subtype}(b_X) \simeq_t \mathrm{Subtype}(b_Y))$.
\end{enumerate}
Then $\mathrm{IsEmpty}(X \simeq_t Y)$.
\end{theorem}

\begin{proof}[Proof sketch]
Any homeomorphism $e : X \simeq_t Y$ preserves $\mathrm{ChartableR2}$
(Theorem~\ref{thm:chartable-invar}), hence maps $b_X$ to $b_Y$ pointwise.
By $e.\mathrm{subtype}$, we get a homeomorphism
$\mathrm{Subtype}(b_X) \simeq_t \mathrm{Subtype}(b_Y)$, contradicting~(iii).
\end{proof}

% This means: to prove IsEmpty (X ≃ₜ Y) for any NEW pair, supply boundary
% predicates + ChartableR2 biconditionals + boundary incompatibility.
% No bridge re-proof needed. M-FINAL is a one-line corollary.
% Klein bottle vs torus, etc. would follow the same pattern.

\subsection{Half-Space Obstructions}
\label{sec:halfspace}

% The engine that converts "boundary point with H2-model neighborhood"
% into "not ChartableR2".

\begin{theorem}[\leanref{HalfPlaneVsPlane.lean}]
\label{thm:halfplane}
(\leanref{euclideanHalfSpace2\_not\_homeomorphic\_euclidean2})
$\mathrm{IsEmpty}(H^2 \simeq_t \mathbb{R}^2)$.
Proof via punctured half-plane (star-convex $\Rightarrow$ simply connected) vs.\ punctured plane (not simply connected, from $\exp$ covering).
\end{theorem}

\begin{theorem}[\leanref{HalfSpaceNeighborVsPlane.lean}]
\label{thm:halfspace-local}
(\leanref{isEmpty\_homeomorph\_subtype\_openH2\_zero\_nbhd\_subtype\_openPlane})
For open $V \subseteq H^2$ with $0 \in V$ and open $O \subseteq \mathbb{R}^2$, there is no homeomorphism $V \cong_t O$.
\end{theorem}

% Supporting: HalfLineVsLine (1D), PuncturedPlaneNotSimplyConnected.

\subsection{Cylinder C4: Boundary Biconditional}

\begin{theorem}[\leanref{CylinderChartableBoundary.lean}]
\label{thm:cylinder-c4}
(\leanref{closedCylinder\_boundaryUnion\_iff\_not\_chartableR2})
For all $x \in \mathrm{ClosedCylinder}$:
\[
  x \in \partial_{\mathrm{Cyl}} \;\iff\; \neg\,\mathrm{ChartableR2}(x).
\]
\end{theorem}

% Proof: interior points get product charts ≃ R²; boundary points modeled
% via homeomorph_closedCylinder_model_to_H2 + not_chartableR2_of_isOpenEmbedding_H2.
% Lean: chartableR2_closedCylinder_of_Ioo (interior), not_chartableR2 lemmas (boundary).

% ============================================================================
% 5. THE SEAM-CHART PROBLEM
% ============================================================================
\section{Interior Charts on the M\"obius Strip: The Seam Problem}
\label{sec:seam}

% ~5 pages. This is the hardest part of the proof and the most original
% proof technique.

\subsection{Why the Seam Is Hard}
\label{sec:seam-why}

% Fundamental domain [0,1]^2. Three zones for interior-height points (0 < t < 1):
%   (i)   Open cell: 0 < x < 1. Quotient map is injective here; standard
%         open-box chart trivially lifts. [Easy]
%   (ii)  Seam away from equator: x ∈ {0,1}, t ≠ ½. Trigonometric coordinate
%         flattening is injective when δ < |t₀ - ½|. [Medium]
%   (iii) Equator point: (0, ½) ~ (1, ½). The trig chart degenerates (the
%         t₀ = ½ case breaks the δ < |t₀ - ½| window). A different chart is
%         needed. [Hard — novel technique]

\subsection{Open-Cell Charts}

\begin{theorem}[\leanref{MobiusChartableBoundary.lean}]
\label{thm:open-cell}
(\leanref{chartableR2\_mobiusQuotientMk\_of\_Ioo\_square})
If $p \in (0,1)^2$ (strict interior of the fundamental domain), then $\mathrm{ChartableR2}(\llbracket p \rrbracket)$.
\end{theorem}

% Proof: quotient map injective on Ioo square (mobiusClass_eq_singleton_of_Ioo_fx);
% image of open set under injective continuous map from compact Hausdorff = open;
% homeomorphism with R² via mobiusPlaneBoxHomeomorph.

\subsection{Seam Trig Charts (Away from Equator)}

\begin{definition}[\leanref{MobiusSeamChart.lean}, \leanref{MobiusSeamChartable.lean}]
The \emph{saturated seam patch} $S(t_0, \varepsilon, \delta)$ is the union of the left and right $\varepsilon$-thickened seam strips at height window $(t_0 - \delta, t_0 + \delta)$, saturated under $\sim$.
\end{definition}

\begin{theorem}[\leanref{MobiusSeamTrigInject.lean}]
\label{thm:trig-inject}
(\leanref{onPatch\_mobiusFDTrigCoord\_eq\_iff\_mobiusRel0})
When $\delta < |t_0 - \tfrac{1}{2}|$, the trigonometric FD coordinate map $\mathrm{mobiusFDTrigCoord}$ satisfies: $\mathrm{mobiusFDTrigCoord}(p) = \mathrm{mobiusFDTrigCoord}(q)$ on the saturated patch if and only if $p \sim q$.
\end{theorem}

\begin{theorem}[\leanref{MobiusSeamChartableR2.lean}]
\label{thm:seam-trig}
(\leanref{chartableR2\_mobiusQuotientMk\_of\_mem\_mobiusSeamSaturatedPatch})
For $t_0 \neq \tfrac{1}{2}$ and appropriate $\varepsilon, \delta$, every point in $\pi(S(t_0, \varepsilon, \delta))$ is $\mathrm{ChartableR2}$.
\end{theorem}

% Proof sketch: injectivity (trig) + open image (B3) + quotient homeomorphism
% with open subset of R².

\subsection{The Equator Chart: Piecewise Unfolding}
\label{sec:equator}

% THIS IS THE NOVEL PROOF TECHNIQUE.

\begin{definition}[\leanref{MobiusSeamChartableR2.lean}]
\label{def:unfold}
The \emph{fundamental-domain unfolding map} $\mathrm{unfoldFD} : [0,1]^2 \to \mathbb{R}^2$ is:
\[
  \mathrm{unfoldFD}(x, t) :=
  \begin{cases}
    \mathrm{mobiusPlaneCoord}(x, t) & \text{if } x \leq \tfrac{1}{2}, \\
    \mathrm{mobiusPlaneTrigSeamPartnerSub}(\mathrm{mobiusPlaneCoord}(x, t)) & \text{if } x > \tfrac{1}{2}.
  \end{cases}
\]
\end{definition}

\begin{lemma}[\leanref{MobiusSeamChartableR2.lean}]
$\mathrm{unfoldFD}$ is $\sim$-invariant: if $p \sim q$ then $\mathrm{unfoldFD}(p) = \mathrm{unfoldFD}(q)$.
\end{lemma}

% Proof: both formulas agree on (0,t) ~ (1,1-t) because
% mobiusPlaneTrigSeamPartnerSub maps the right-side plane coordinate
% to the left-side plane coordinate of the partner.

\begin{theorem}[\leanref{MobiusSeamChartableR2.lean}]
\label{thm:equator}
(\leanref{chartableR2\_mobiusQuotientMk\_of\_left\_seam\_half})
$\mathrm{ChartableR2}(\llbracket (0, \tfrac{1}{2}) \rrbracket)$.
\end{theorem}

% Proof sketch:
%   1. unfoldFD lifts to mobiusStripUnfold : MobiusStrip → R² (quotient descent).
%   2. Let V = π''{ q ∈ S | unfoldFD(q) ∈ Ball(v, r) } where S is the
%      saturated patch at (½, ¼, ¼) and v = unfoldFD(0, ½).
%   3. V is open (S open + unfoldFD continuous on each half + preimage of open ball).
%   4. V is saturated (unfoldFD is ~-invariant).
%   5. mobiusStripUnfold restricted to V is injective (case split: both in left
%      patch, both in right patch, one in each → unfoldFD values match iff ~ related).
%   6. Image mobiusStripUnfold''V is open in R² (saturated open quotient image +
%      unfoldFD maps onto an open rectangle).
%   7. Build homeomorphism V ≃ₜ (open subset of R²) ≃ₜ R².
%
% Key helper: continuous_quotientLift_unfold_on_saturatedImage (quotient descent
% for the piecewise section; continuity from pasting lemma on left/right halves).

\subsection{Assembly: All Interior Points Are ChartableR2}

\begin{theorem}[\leanref{MobiusSeamChartableR2.lean}]
\label{thm:interior-height}
(\leanref{chartableR2\_mobiusQuotientMk\_of\_interior\_height})
For every $p \in [0,1]^2$ with $0 < p_2 < 1$:
\[
  \mathrm{ChartableR2}(\llbracket p \rrbracket).
\]
\end{theorem}

% Proof: case split on p₁:
%   p₁ = 0 → left seam (§5.3 for t = ½, §5.2 for t ≠ ½)
%   p₁ = 1 → right seam (symmetric)
%   0 < p₁ < 1 → open-cell chart (§5.1)

\subsection{M\"obius C4: Boundary Biconditional}

\begin{theorem}[\leanref{MobiusChartableBoundary.lean}]
\label{thm:mobius-c4}
(\leanref{mobiusStripBoundary\_iff\_not\_chartableR2\_of\_forall\_off\_edge\_chartable})
For all $x \in \mathrm{MobiusStrip}$:
\[
  x \in \partial_{\mathrm{Mob}} \;\iff\; \neg\,\mathrm{ChartableR2}(x).
\]
\end{theorem}

% Forward: not_chartableR2_of_mem_mobiusStripBoundary (H2 collar models +
% sequential stability at corners).
% Reverse: from chartableR2_mobiusQuotientMk_of_interior_height.

% ============================================================================
% 6. M-FINAL: THE NON-HOMEOMORPHISM
% ============================================================================
\section{M-FINAL: The M\"obius Strip Is Not Homeomorphic to the Closed Cylinder}
\label{sec:mfinal}

% ~2 pages.

\begin{theorem}[\leanref{MobiusSeamChartableR2.lean}, \leanref{ChartableR2Bridge.lean}]
\label{thm:mfinal}
(\leanref{mobiusStrip\_not\_homeomorphic\_closedCylinder})
\[
  \mathrm{IsEmpty}\!\left(\mathrm{MobiusStrip} \simeq_t \mathrm{ClosedCylinder}\right).
\]
\end{theorem}

\begin{proof}[Proof sketch]
Instantiate Theorem~\ref{thm:generic} with $b_X = \partial_{\mathrm{Mob}}$,
$b_Y = \partial_{\mathrm{Cyl}}$, the biconditionals
(Theorems~\ref{thm:cylinder-c4} and~\ref{thm:mobius-c4}), and the boundary
incompatibility (Theorem~\ref{thm:boundary-contrast}).
\end{proof}

% The Lean proof chain:
%   1. not_homeomorphic_of_chartableR2_boundary_contrast (generic engine)
%   2. mobiusStrip_not_homeomorphic_closedCylinder_of_chartable_boundary (instantiates generic)
%   3. mobiusStrip_not_homeomorphic_closedCylinder_of_mobius_boundary_chartable (cylinder C4 discharged)
%   4. mobiusStrip_not_homeomorphic_closedCylinder_of_forall_off_edge_chartable (Möbius C4 reduced)
%   5. mobiusStrip_not_homeomorphic_closedCylinder (instantiates with interior_height)

% ============================================================================
% 7. THE EXTENDED MASTER THEOREM AND INSTANTIATIONS
% ============================================================================
\section{The Extended Master Theorem and Instantiations}
\label{sec:extended}

\subsection{Formal Statement}

\begin{theorem}[\leanref{Main.lean}]
\label{thm:extended}
(\leanref{representational\_regress\_master\_extended})
For any representational system $R$:
\[
  \mathrm{RepresentationalRegressMasterExtended}(R)
\]
where this abbreviation expands to:
\begin{enumerate}[label=(\roman*)]
\item $\mathrm{representational\_regress\_master\_claim}(R)$ (regress + wall + $T_2$ transport),
\item $\mathrm{IsEmpty}(H^1 \simeq_t \mathbb{R}^1)$ \quad (1D half-space vs.\ line; \leanref{representational\_regress\_topology\_halfLineModel}),
\item $\mathrm{IsEmpty}(\mathrm{MobiusStrip} \simeq_t \mathrm{ClosedCylinder})$ (M-FINAL).
\end{enumerate}
\end{theorem}

\noindent
This \emph{extends} \leanref{representational\_regress\_master}; it does not replace it.

\subsection{Bundled boundary models}

\begin{definition}[\leanref{ChartableR2Neighbor.lean}]
A \emph{\texttt{ChartableR2} boundary model} on $X$ is a pair $(b_X,\beta_X)$ where
$b_X \subseteq X$ and $\beta_X$ witnesses
$\forall x,\; x \in b_X \iff \neg\,\mathrm{ChartableR2}(x)$.
\end{definition}

\noindent
Lean provides \leanref{ChartableR2BoundaryModel} and
\leanref{ChartableR2BoundaryModel.not\_homeomorphic}. Concrete instances
\leanref{mobiusStripChartableR2BoundaryModel} and
\leanref{closedCylinderChartableR2BoundaryModel} sit in
\leanref{ChartableR2ConcreteBoundaryModels}; M-FINAL is recovered as
\leanref{mobiusStrip\_closedCylinder\_boundary\_models\_not\_homeomorphic}.

\subsection{Instantiations}
\label{sec:instantiations}

Theorems below are either the same mathematics under paper-facing names or
\emph{conditional} schemes: no alternative axioms and no quantum formalism.

\subsubsection{Consciousness and self-presence (conditional)}

Premise \textbf{P2} (not in Lean): a target phenomenon of interest---call it
\emph{conscious self-presence}---is modeled topologically by a space globally
like $\mathrm{MobiusStrip}$ (non-orientable side collapse, one boundary band in
our witness pair) rather than like $\mathrm{ClosedCylinder}$ (two-sided cut).

\medskip
\noindent\textbf{Formal content (unconditional).}
Theorem~\ref{thm:mfinal}: cylinder and M\"obius strip are not homeomorphic;
hence no homeomorphism identifies a ``cylinder-class'' model with a
``M\"obius-class'' model \emph{if} boundary sensors and subspace data match the
proved pattern.

\medskip
\noindent\textbf{Conditional gloss.}
If one's philosophical story places representational iteration on the cylinder
side of that contrast and self-presence on the M\"obius side, then the gap is
\emph{topological}, not a deficit of compute, recursion depth, or training
data. Rejecting P2 is coherent; Lean does not adjudicate P2.

\subsubsection{Symbol grounding and semantic regress}

\leanref{SymbolSystem} is definitionally \leanref{RepresentationalSystem}; lemmas
\leanref{symbolGrounding\_metaRegressArrow\_injective} (etc.) restate that
distinct meta-levels of ``interpretation'' are distinct morphisms. This makes
the grounding--regress connection citeable without a parallel definition tree.

\subsubsection{Diagonal arguments and self-interpretation}

\leanref{LawvereRegressBridge} packages:
\leanref{no\_universal\_parametric\_unary} /
\leanref{no\_universal\_parametric\_unary\_nat} (no family $s : A \to A \to B$
listing all $A \to B$ when $B$ has a fixed-point-free endomap---Cor.~to Lawvere
in \texttt{Type}), and
\leanref{lawvere\_fixed\_point\_stays\_representational}: iterates of
\texttt{represent} never occupy the \texttt{obj} branch of \texttt{OntologicalSlot}.

\subsubsection{Quantum measurement chains (schematic)}

\leanref{QuantumObserverChainHypothesis} bundles exactly the hypotheses needed to
apply Theorem~\ref{thm:generic}: a topological space of ``observers + readouts,''
a boundary set with the \texttt{ChartableR2} biconditional, and subtype
non-homeomorphism to $\mathrm{mobiusStripBoundary}$. Then
\leanref{QuantumObserverChainHypothesis.not\_homeomorphic\_mobiusStrip} yields
$\mathrm{IsEmpty}(\mathrm{Obs} \simeq_t \mathrm{MobiusStrip})$.
\textbf{No physical axioms:} instantiating this requires a future formal model of
the state space and boundary sensor.

\subsection{What this means for AI and scaling}

More parameters, longer context, richer self-models, and deeper recurrent stacks
do not constitute a homeomorphism of a representational substrate onto a
different global topology: homeomorphism invariance is the mathematical content
behind ``scale does not change the invariant.''

\subsection{Relationship to the Trans-Computational Processing (TP) argument}

The TP line uses computability (Halting, G\"odel-style diagonalization). Here the
independent spine is topological (\texttt{ChartableR2}, boundary components). Both
support a \emph{limitative} moral about self-contained representation; the lemmas
are not interchangeable but the narrative arcs rhyme.

% ============================================================================
% 8. HONEST LIMITS
% ============================================================================
\section{Honest Limits}
\label{sec:limits}

% ~2 pages. Critical for credibility.

\begin{enumerate}
\item \textbf{Hypothesis, not derivation.} $\mathrm{iter\_injective}$ is bundled, not proved from minimal data. Not every $f : A \to A$ in every category satisfies it. The regress is conditional on this non-triviality hypothesis.

\item \textbf{Lawvere in \texttt{Type}, not abstract CCC.} The formalization uses $\mathrm{Type}\ u$ and $\mathrm{MonoidalClosed}(\mathrm{Type}\ u)$, not a general closed category over arbitrary $C$. A fully abstract CCC version remains future work.

\item \textbf{\texttt{OntologicalSlot} is definitional, not metaphysical.} The disjointness $\mathrm{endo} \neq \mathrm{obj}$ is trivially true by constructor discrimination. This is the point, but the triviality means the result is \emph{semantic} (about the structure of type theory) rather than a deep theorem.

\item \textbf{Models are chosen, not unique.} The closed cylinder and M\"obius strip are \emph{specific} topological models chosen for clarity and tractability. A different philosophical account might use different geometric witnesses (e.g., Klein bottle, different quotient). The argument works for any two models separated by a topological invariant; these two are the cleanest.

\item \textbf{P2 is philosophical.} The premise ``consciousness requires non-orientable topology'' is not a theorem. The formalization proves the conditional: \emph{if P2, then}~\ldots{}. Denying P2 is coherent; the formalization does not settle whether one should accept it.

\item \textbf{No $\mathrm{IsManifold}$ polish.} The M\"obius strip is not given a full Mathlib $\mathrm{IsManifold}$ instance. The proof uses $\mathrm{ChartableR2}$ (open neighborhoods homeomorphic to $\mathbb{R}^2$) as a self-contained predicate, sidestepping the need for a full manifold-with-boundary atlas. This is sufficient for the non-homeomorphism but less elegant than a full atlas approach.

\item \textbf{Route W is incomplete.} An alternative proof via circle-winding / doubling obstructions (Steps~1--4 in \leanref{CylinderMobiusNonhomeo.lean}) is partially formalized. Step~5 (transport under a hypothetical global homeomorphism) is not proved. The flagship proof via $\mathrm{ChartableR2}$ is complete and does not depend on Route~W.

\item \textbf{\texttt{QuantumObserverChainHypothesis} is not physics.} It records a
  \emph{pattern} of assumptions. No Hilbert space, operators, or measurement
  axioms appear; any physical reading requires proving those assumptions in a
  separate formalization.
\end{enumerate}

% ============================================================================
% 9. RELATED WORK
% ============================================================================
\section{Related Work}
\label{sec:related}

% ~2 pages.

\subsection{Formalized Topology in Proof Assistants}
% Mathlib's manifold library (Sébastien Gouëzel, Heather Macbeth, et al.).
% sphere eversion (Patrick Massot, Oliver Nash, et al.).
% Formal proofs of topological facts in Coq, Isabelle.
% Our contribution: application to a philosophical argument, not pure math.

\subsection{Lawvere's Theorem in Proof Assistants}
% Prior Coq/Agda formalizations of Lawvere's fixed-point theorem.
% Our addition: OntologicalSlot construction (novel).

\subsection{Mathematical Arguments About Consciousness}
% Penrose (Gödel arguments); Tononi (IIT, information geometry);
% Hofstadter (strange loops — orientable, per our analysis);
% Chalmers (hard problem); TP paper (Spivack 2025, computability).
% Our contribution: first machine-checked topological argument.

\subsection{Topology and Philosophy}
% Existing uses of topology in philosophy of mind (rare and informal).
% Our contribution: specific invariant (boundary component count), proved,
% not metaphorical.

% ============================================================================
% 10. CONCLUSION
% ============================================================================
\section{Conclusion}
\label{sec:conclusion}

We presented a machine-checked library for representational regress, Lawvere
packaging, a generic $\mathrm{ChartableR2}$ boundary obstruction, and the
flagship non-homeomorphism $\mathrm{MobiusStrip} \not\simeq_t
\mathrm{ClosedCylinder}$. The mathematics stands without phenomenology; the
same lemmas support cited instantiations (consciousness under P2, symbol
grounding, diagonal limitations, and a schematic measurement-chain interface)
without widening proof obligations.

\medskip
\noindent\textbf{Future work:} full $\mathrm{IsManifold}$ polish for
$\mathrm{MobiusStrip}$; Route~W completion; abstract CCC Lawvere; deeper links
to other formal theories where boundary data admit $\mathrm{ChartableR2}$
sensors.

% ============================================================================
% APPENDICES
% ============================================================================
\appendix

\section{Lean Module Map}
\label{app:modules}

% Table: one row per module, columns: Path | Role | Key theorems.
% Directly from REPRESENTATIONAL_REGRESS_FORMALIZATION_MAP.md.

\begin{table}[h]
\centering
\small
\begin{tabular}{@{}lp{8cm}@{}}
\toprule
\textbf{Module} & \textbf{Role} \\
\midrule
\leanref{Basic} & \leanref{RepresentationalSystem}, iterates, \leanref{OntologicalSlot} \\
\leanref{Regress} & Non-termination, unbounded levels \\
\leanref{FixedPoints} & \leanref{EndoVsPoint}, Lawvere wall non-dissolution \\
\leanref{Lawvere} & Fixed-point theorem in \texttt{Type} \\
\leanref{LawvereCCCType} & \texttt{MonoidalClosed} packaging \\
\leanref{Orientability} & $T_2$ homeomorphism invariance \\
\leanref{CylinderMobius} & Fundamental domain, quotient, compact Hausdorff, boundary band \\
\leanref{CylinderBoundary} & Closed cylinder boundary: two disjoint connected faces \\
\leanref{CylinderChartableBoundary} & Cylinder C4 biconditional \\
\leanref{ChartableR2Neighbor} & \leanref{ChartableR2}, \leanref{ChartableR2BoundaryModel}; openness / stability \\
\leanref{ChartableR2Bridge} & \leanref{Homeomorph.chartableR2\_iff}; \leanref{not\_homeomorphic\_of\_chartableR2\_boundary\_contrast}; \leanref{ChartableR2BoundaryModel.not\_homeomorphic} \\
\leanref{ChartableR2ConcreteBoundaryModels} & M\"obius/cylinder boundary models; conditional observer chain; \leanref{QuantumObserverChainHypothesis} \\
\leanref{SymbolGrounding} & \leanref{SymbolSystem}; \leanref{symbolGrounding\_*} \\
\leanref{LawvereRegressBridge} & \leanref{no\_universal\_parametric\_unary\_nat}; \leanref{lawvere\_fixed\_point\_stays\_representational} \\
\leanref{ChartableR2Models} & Interior charts; cylinder interior \leanref{ChartableR2} \\
\leanref{MobiusChartableBoundary} & Open-cell charts; boundary $\neg$\leanref{ChartableR2}; M\"obius C4 packaging \\
\leanref{MobiusSeamChartable} & Saturated seam patches (open, saturated) \\
\leanref{MobiusSeamChart} & IFT-based local model; trig coordinates \\
\leanref{MobiusSeamTrigInject} & Trig injectivity on quotient windows \\
\leanref{MobiusSeamChartableR2} & Seam charts (trig + unfold); interior-height; \textbf{M-FINAL} \\
\leanref{MobiusCylinderBoundaryContrast} & Boundary subspaces not homeomorphic \\
\leanref{CylinderMobiusNonhomeo} & Route W (optional; Steps 1--4) \\
\leanref{HalfLineVsLine} & 1D half-space $\neq$ line \\
\leanref{HalfPlaneVsPlane} & 2D half-space $\neq$ plane \\
\leanref{HalfSpaceNeighborVsPlane} & Local half-space $\neq$ open plane patch \\
\leanref{PuncturedPlaneNotSimplyConnected} & $\mathbb{C} \setminus \{0\}$ not simply connected \\
\leanref{OpenCompactChartObstruction} & Compact $\mathbb{R}^2$ patch $\neq$ $\mathbb{R}^2$ \\
\leanref{Main} & Master theorem + extended bundle \\
\bottomrule
\end{tabular}
\end{table}

\section{Key Theorem Index}
\label{app:theorems}

% Compact table of the ~20 most important theorems with Lean names.
% Columns: Lean name | Informal statement | Module.
% From THEOREM_INVENTORY.md (curated selection).

\begin{table}[h]
\centering
\small
\begin{tabular}{@{}lp{6.5cm}l@{}}
\toprule
\textbf{Lean name} & \textbf{Statement} & \textbf{Module} \\
\midrule
\leanref{regress\_iterates\_unbounded} & $\forall B, \exists L > B$, distinct iterate & \leanref{Regress} \\
\leanref{lawvere\_fixed\_point\_Type} & Universal enumerator $\Rightarrow$ every endo has fixed point & \leanref{Lawvere} \\
\leanref{lawvere\_wall\_is\_not\_dissolution} & $\mathrm{endo}(\mathrm{represent}) \neq \mathrm{obj}(\mathrm{fp})$ & \leanref{FixedPoints} \\
\leanref{instT2SpaceMobiusStrip} & M\"obius strip is $T_2$ & \leanref{CylinderMobius} \\
\leanref{isConnected\_mobiusStripBoundary} & Boundary band is connected & \leanref{CylinderMobius} \\
\leanref{closedCylinder\_boundary\_eq\_union} & Cylinder boundary = 2 disjoint faces & \leanref{CylinderBoundary} \\
\leanref{euclideanHalfSpace2\_not\_homeomorphic\_euclidean2} & $H^2 \not\cong_t \mathbb{R}^2$ & \leanref{HalfPlaneVsPlane} \\
\leanref{not\_homeomorphic\_of\_chartableR2\_boundary\_contrast} & Generic: incompatible boundaries $\Rightarrow$ $\neg$ homeo & \leanref{ChartableR2Bridge} \\
\leanref{closedCylinder\_boundaryUnion\_iff\_not\_chartableR2} & Cylinder C4 biconditional & \leanref{CylinderChartableBoundary} \\
\leanref{not\_chartableR2\_of\_mem\_mobiusStripBoundary} & M\"obius boundary $\Rightarrow$ $\neg$\leanref{ChartableR2} & \leanref{MobiusChartableBdy} \\
\leanref{chartableR2\_mobiusQuotientMk\_of\_interior\_height} & Interior height $\Rightarrow$ \leanref{ChartableR2} & \leanref{MobiusSeamChartableR2} \\
\leanref{mobiusStripBoundary\_not\_homeomorphic\ldots} & Boundary subtypes not homeo & \leanref{MobiusCylBdyContrast} \\
\leanref{mobiusStrip\_not\_homeomorphic\_closedCylinder} & \textbf{M-FINAL} & \leanref{MobiusSeamChartableR2} \\
\leanref{representational\_regress\_master} & Core master conjunction & \leanref{Main} \\
\leanref{representational\_regress\_master\_extended} & Extended bundle (+ topology) & \leanref{Main} \\
\leanref{QuantumObserverChainHypothesis.not\_homeomorphic\_mobiusStrip} & Conditional: obs.\ chain $\not\simeq_t$ M\"obius & \leanref{ChartableR2ConcreteBoundaryModels} \\
\leanref{no\_universal\_parametric\_unary\_nat} & No universal $s : A \!\to\! A \!\to\! \mathbb{N}$ for all $g : A \!\to\! \mathbb{N}$ & \leanref{LawvereRegressBridge} \\
\bottomrule
\end{tabular}
\end{table}

\section{Reproduction}
\label{app:repro}

% Build instructions, toolchain versions, verification commands.
% From ARTIFACT.md.

\begin{verbatim}
cd representational-regress-lean
lake update
lake exe cache get
lake build RepresentationalRegress
# Verify: 0 sorry
rg "sorry" RepresentationalRegress/
\end{verbatim}

\noindent
Toolchain: \texttt{leanprover/lean4:v4.29.0-rc6}. Mathlib: \texttt{v4.29.0-rc6}.

\end{document}
